\begin{abstract}
This paper proposes Information Field Geometry Theory (IFGT). IFGT is a framework for describing information not as a substance or a description of states, but as a \emph{quasi-closure}: a residual structure that persists when difference does not fully close in the process of approaching closure.

The theory builds on Vortical Enclosure Dynamics (VED), which describes the generation of physical structure from difference, gradient, flow, and closure. Within the same generative process described by VED, IFGT treats the persisting unclosed structure as information. While VED describes physical persistence realized through the closure degree $C$, IFGT addresses the complementary side---what fails to fully close---as information. In this framing, information is not opposed to physical structure but is a complementary aspect of the same generative structure.

IFGT introduces five basic variables: information density $I$, information flow $J$, information potential $\Phi$, informational time $\tau$, and the generation term $\sigma$. The relation between $I$ and $C$ is given by $I = f(1 - C)$, and $\sigma$ is defined as the mismatch between the current difference chain and the existing trace structure. In the final chapter, these components are integrated into a unified VED$\times$IFGT equation, yielding a single evolution equation for the closure-degree field $C(x,\tau)$.

Within this framework, information is not a static object of description but a dynamic process that participates in structural formation through the persistence of difference and the constraint of flow. IFGT thus provides a foundational theory for describing the general structure that may encompass informational behavior observed in biological and computational systems.
\end{abstract}
